\section*{Bloque D. Ampliación: Teoría y Aplicaciones Prácticas (70\% / 30\%)}

\noindent\textit{Preguntas 91 a 180. Ampliación del banco oficial con 63 preguntas teóricas y conceptuales (91 a 153) y 27 aplicaciones cuantitativas prácticas (154 a 180).}

\subsection*{Parte I. Preguntas teóricas y conceptuales (91 a 153)}

\q{¿Por qué el supuesto de homogeneidad del producto es indispensable para que las empresas sean precio-aceptantes en competencia perfecta?}
\begin{opts}
    \item Porque garantiza que los costos fijos medios de todas las empresas sean idénticos en el corto plazo.
    \item Porque elimina la diferenciación y hace que los bienes de distintas firmas sean sustitutos perfectos.
    \item Porque asegura que la demanda agregada del mercado sea perfectamente elástica en cualquier nivel de precios.
    \item Porque impide que existan economías externas pecuniarias cuando la industria se expande en el largo plazo.
\end{opts}

\q{En un mercado perfectamente competitivo en equilibrio, el valor del Índice de Lerner $L = \frac{P - CMg}{P}$ es igual a:}
\begin{opts}
    \item Uno, debido a que la elasticidad precio de la demanda del mercado es unitaria en el óptimo.
    \item Infinito, dado que la curva de demanda individual de la empresa tiene pendiente horizontal.
    \item Cero, porque la firma iguala el precio al costo marginal y carece de poder de mercado.
    \item El inverso negativo de la elasticidad precio de la demanda agregada de toda la industria.
\end{opts}

\q{¿Cómo se concilia que la curva de demanda del mercado competitivo tenga pendiente negativa mientras la demanda que enfrenta cada firma individual es horizontal?}
\begin{opts}
    \item Cada firma aporta una fracción ínfima de la oferta total y sus cambios de volumen no alteran el precio.
    \item Los consumidores tienen preferencias cuasilineales idénticas que anulan el efecto ingreso en la demanda individual de cada firma.
    \item La agregación vertical de curvas perfectamente elásticas genera automáticamente una curva de mercado con pendiente negativa.
    \item Las empresas coordinan sus volúmenes de producción para mantener constante el precio de mercado ante cambios en la demanda.
\end{opts}

\q{Si en un mercado con muchos oferentes y bien homogéneo se viola el supuesto de información perfecta sobre los precios, la consecuencia más directa es que:}
\begin{opts}
    \item El precio de equilibrio cae por debajo del mínimo del costo variable medio de todas las empresas.
    \item Todas las empresas abandonan inmediatamente el mercado incluso obteniendo beneficios contables positivos.
    \item La curva de oferta de la industria a largo plazo se vuelve necesariamente de pendiente negativa.
    \item Puede surgir dispersión de precios y cierto poder de mercado local debido a los costos de búsqueda.
\end{opts}

\q{En el modelo de competencia perfecta, el supuesto de libre movilidad de factores y ausencia de barreras a la entrada cumple principalmente la función de:}
\begin{opts}
    \item Impedir que las empresas sufran pérdidas económicas durante el corto plazo ante choques adversos.
    \item Asegurar que los beneficios económicos extraordinarios se disipen en el largo plazo vía entrada y salida.
    \item Garantizar que todas las empresas operen con rendimientos constantes a escala en todo nivel de producción.
    \item Mantener inalterados los precios de los insumos primarios cuando se expande la producción de la industria.
\end{opts}

\q{Para una empresa precio-aceptante, la función de ingreso total $IT(q) = P \cdot q$ es lineal desde el origen, lo cual implica que:}
\begin{opts}
    \item El ingreso marginal es siempre estrictamente menor que el ingreso medio para todo nivel positivo de $q$.
    \item El ingreso medio disminuye conforme aumenta la producción, mientras el ingreso marginal permanece constante.
    \item El ingreso marginal, el ingreso medio y el precio de mercado coinciden exactamente para cualquier $q > 0$.
    \item La elasticidad ingreso de la demanda del bien producido por la empresa es estrictamente igual a la unidad.
\end{opts}

\q{En un mercado competitivo ideal, la existencia de costos de transacción nulos garantiza que:}
\begin{opts}
    \item No exista brecha entre el precio efectivamente pagado por el comprador y el precio recibido por el vendedor.
    \item El costo marginal de producción de todas las empresas de la industria sea exactamente igual a cero.
    \item Las empresas puedan diferenciar sus productos sin incurrir en gastos adicionales de publicidad o diseño.
    \item El gobierno pueda recaudar impuestos específicos sin generar pérdida irrecuperable de eficiencia en el mercado.
\end{opts}

\q{Según la teoría de mercados disputables (contestable markets), un mercado con pocas empresas puede aproximar el resultado competitivo ($P = CMe$) incluso sin atomicidad si:}
\begin{opts}
    \item El regulador fija un precio mínimo por encima del costo medio de corto plazo de la firma dominante.
    \item Los consumidores presentan una demanda completamente inelástica frente a incrementos en el precio del bien.
    \item Todas las firmas comparten una tecnología con costos fijos hundidos elevados que impiden la salida rápida.
    \item Existe entrada y salida totalmente libres sin costos hundidos (entrada relámpago hit-and-run).
\end{opts}

\q{Si una empresa competitiva encuentra un nivel de producción $q_0 > 0$ donde $P = CMg(q_0)$, pero en ese punto $\frac{dCMg}{dq} < 0$, entonces $q_0$ corresponde a:}
\begin{opts}
    \item Un máximo global del beneficio económico, siempre que el precio supere al costo fijo medio.
    \item Un mínimo local del beneficio económico, por lo que la firma no debe producir esa cantidad.
    \item El punto de cierre de corto plazo donde el costo marginal intercepta al mínimo del costo variable medio.
    \item Un equilibrio de largo plazo con beneficio económico nulo y eficiencia productiva plena.
\end{opts}

\q{En el corto plazo, si el precio de mercado satisface $\min CVMe < P < \min CMe$, la decisión óptima de la empresa competitiva es:}
\begin{opts}
    \item Cerrar inmediatamente porque obtiene pérdidas económicas y no logra cubrir sus costos totales medios.
    \item Producir la cantidad donde $P = CMe$ para asegurar que el beneficio contable sea exactamente cero.
    \item Seguir produciendo donde $P = CMg$, pues cubre sus costos variables y amortigua parte del costo fijo.
    \item Aumentar la producción hasta igualar el costo marginal con el mínimo del costo total medio.
\end{opts}

\q{¿Cómo se modifica la regla de cierre en el corto plazo si una fracción del costo fijo es evitable (quasi-fijo) al suspender operaciones?}
\begin{opts}
    \item La empresa cierra si el ingreso total no cubre la suma del costo variable más el costo fijo evitable.
    \item La empresa ignora tanto los costos fijos evitables como los hundidos y solo compara el precio con el $CVMe$.
    \item La empresa produce siempre que el precio supere únicamente al costo fijo hundido dividido por la cantidad.
    \item La curva de oferta de corto plazo se vuelve horizontal al nivel del costo fijo total de la firma.
\end{opts}

\q{Si el Estado impone un impuesto de suma alzada (licencia anual fija independiente de $q$) a las empresas de una industria competitiva que ya operan en el corto plazo, entonces:}
\begin{opts}
    \item La curva de costo marginal se desplaza hacia arriba y cada empresa reduce su producción óptima.
    \item El precio de equilibrio de corto plazo aumenta inmediatamente en la misma proporción que el impuesto.
    \item La curva de oferta de corto plazo de la industria se contrae porque sube el mínimo del costo variable medio.
    \item La producción óptima y el precio de equilibrio no cambian en el corto plazo, pero caen los beneficios.
\end{opts}

\q{¿Bajo qué circunstancia la curva de oferta de corto plazo de la industria competitiva es más inelástica que la suma horizontal simple de las curvas de $CMg$ individuales?}
\begin{opts}
    \item Cuando todas las empresas de la industria poseen exactamente la misma función de costos cuadrática.
    \item Cuando la expansión simultánea de todas las empresas eleva el precio de los insumos variables.
    \item Cuando el número de empresas en la industria se mantiene fijo durante todo el período de análisis.
    \item Cuando la demanda de mercado es perfectamente elástica respecto del precio del bien final.
\end{opts}

\q{La pendiente positiva de la curva de oferta de corto plazo de una empresa competitiva se explica directamente por:}
\begin{opts}
    \item La existencia de economías externas pecuniarias en el mercado de factores productivos.
    \item El incremento del costo fijo medio conforme la empresa amplía su escala de planta.
    \item La ley de los rendimientos marginales decrecientes del factor variable en el corto plazo.
    \item La concavidad de la curva de demanda agregada que enfrentan las empresas de la industria.
\end{opts}

\q{Si se incrementa el precio de una materia prima utilizada en proporción directa a la producción de cada unidad, en el corto plazo competitivo:}
\begin{opts}
    \item Aumenta el $CMg$ y el punto de cierre, contrayendo la oferta industrial y elevando el precio de equilibrio.
    \item Solo se eleva el costo fijo medio sin modificar el nivel de producción ni el precio de mercado.
    \item Disminuye el punto de cierre de cada firma porque los costos hundidos pierden peso relativo en el costo total.
    \item Se desplaza la demanda del mercado hacia la derecha para compensar el mayor costo unitario de las firmas.
\end{opts}

\q{El punto de nivelación (o de beneficio económico nulo) de una empresa perfectamente competitiva ocurre exactamente donde:}
\begin{opts}
    \item El costo marginal corta al costo variable medio en su punto mínimo de corto plazo.
    \item El ingreso total iguala al costo variable total y las pérdidas son iguales al costo fijo.
    \item El ingreso marginal supera al costo medio en una magnitud igual al costo de oportunidad del capital.
    \item El precio de mercado iguala al mínimo del costo total medio, donde $CMg = CMe$.
\end{opts}

\q{Para una empresa competitiva con costos fijos totalmente hundidos, la pérdida económica máxima en la que puede incurrir al tomar su decisión óptima de corto plazo es:}
\begin{opts}
    \item Igual a la suma del costo fijo total más el costo variable medio evaluado en la escala eficiente.
    \item Igual al costo fijo total, pues si $P < \min CVMe$ elige producir cero unidades.
    \item Estrictamente cero, ya que la libre salida garantiza que ninguna empresa pierda dinero en el corto plazo.
    \item Proporcional al cuadrado de la cantidad producida cuando el precio cae por debajo del costo medio.
\end{opts}

\q{En el equilibrio competitivo de largo plazo, las empresas obtienen beneficio económico nulo ($\pi = 0$). ¿Por qué, pese a ello, no abandonan la industria?}
\begin{opts}
    \item Porque los contratos laborales y las barreras legales les impiden salir del mercado en el largo plazo.
    \item Porque el gobierno les otorga transferencias de suma alzada equivalentes a sus costos fijos anuales.
    \item Porque $\pi = 0$ implica cubrir costos explícitos e implícitos, obteniendo beneficio normal.
    \item Porque en el largo plazo el beneficio contable es negativo pero el excedente del productor es infinito.
\end{opts}

\q{A diferencia del corto plazo, un impuesto anual de suma alzada (licencia fija por operar) aplicado a todas las empresas de una industria competitiva con libre entrada provoca en el largo plazo:}
\begin{opts}
    \item Un alza del $\min CMe_{LP}$, mayor precio de equilibrio, mayor escala por firma y menos firmas.
    \item Ningún cambio en el precio de equilibrio ni en el número de empresas, pues el impuesto no afecta al $CMg$.
    \item Una reducción en el precio de largo plazo debido a la entrada masiva de firmas para diluir el costo fijo.
    \item Una caída en la escala mínima eficiente de cada empresa manteniendo constante el tamaño total del mercado.
\end{opts}

\q{En una industria competitiva de costos constantes con empresas idénticas en equilibrio de largo plazo, si la demanda del mercado se duplica, entonces en el nuevo equilibrio:}
\begin{opts}
    \item El precio de largo plazo se duplica y cada empresa produce el doble de su escala mínima eficiente.
    \item El precio de largo plazo sube parcialmente y el número de empresas se mantiene exactamente constante.
    \item Cada empresa reduce su producción a la mitad mientras el precio cae por debajo del costo medio mínimo.
    \item El precio de largo plazo y la escala por firma no varían, mientras el número de firmas se duplica.
\end{opts}

\q{¿Qué origina que la curva de oferta de largo plazo ($S_{LP}$) de una industria perfectamente competitiva tenga pendiente positiva?}
\begin{opts}
    \item La presencia de deseconomías internas a escala que impiden a cualquier empresa operar en su $CMe_{LP}$ mínimo.
    \item El encarecimiento de factores especializados cuando la entrada de nuevas firmas expande la demanda de insumos.
    \item La existencia de poder de mercado colectivo ejercido por las empresas establecidas mediante barreras legales.
    \item La disminución progresiva de la elasticidad precio de la demanda conforme aumenta el ingreso nacional.
\end{opts}

\q{En una industria competitiva de costos decrecientes, un aumento permanente de la demanda conduce en el largo plazo a:}
\begin{opts}
    \item Un precio de equilibrio más alto y una cantidad producida menor que en el equilibrio inicial.
    \item Un precio de equilibrio inalterado y una reducción en el número de empresas activas en el mercado.
    \item Un precio de equilibrio menor y una cantidad mayor, debido a externalidades pecuniarias positivas.
    \item Beneficios económicos positivos permanentes para todas las empresas que ingresan después del choque.
\end{opts}

\q{En una industria competitiva de largo plazo donde coexisten empresas con acceso a insumos de distinta calidad (por ejemplo, tierras más fértiles), ¿qué ocurre con los beneficios?}
\begin{opts}
    \item La firma marginal obtiene beneficio nulo, mientras las inframarginales obtienen renta económica.
    \item Todas las firmas obtienen exactamente el mismo beneficio económico positivo debido a la escasez del insumo.
    \item Las firmas con insumos más productivos fijan un precio menor y expulsan del mercado a la firma marginal.
    \item El precio de equilibrio se fija igual al costo medio mínimo de la firma más eficiente de toda la industria.
\end{opts}

\q{La cuasi-renta en el análisis marshalliano de corto plazo coincide conceptualmente con:}
\begin{opts}
    \item El beneficio económico neto obtenido después de deducir tanto costos variables como costos fijos.
    \item El excedente del consumidor generado por la diferencia entre la disposición a pagar y el precio de mercado.
    \item La pérdida irrecuperable de eficiencia que surge cuando el Estado establece un precio máximo efectivo.
    \item El excedente del productor de corto plazo ($IT - CV = \pi + CF$), que remunera a los factores fijos.
\end{opts}

\q{Suponga que una sola empresa competitiva descubre una innovación no patentable que reduce su $CMe_{LP}$ y luego todas las demás firmas pueden imitarla libremente. En el largo plazo:}
\begin{opts}
    \item La firma innovadora conserva permanentemente beneficios extraordinarios al mantener el precio inicial.
    \item El precio cae al nuevo $\min CMe_{LP}$, se expande $Q$ y el beneficio económico vuelve a ser cero.
    \item El número de empresas aumenta pero el precio de equilibrio permanece igual al costo medio antiguo.
    \item Todas las empresas obtienen beneficios económicos negativos por haber reducido sus costos marginales.
\end{opts}

\q{Si todas las empresas de una industria competitiva tienen exactamente la misma tecnología con rendimientos constantes a escala (costo marginal y medio constantes $c$) y precios de factores fijos, entonces en equilibrio de largo plazo:}
\begin{opts}
    \item El número de empresas y el tamaño individual quedan unívocamente determinados por el nivel de demanda.
    \item El precio de equilibrio depende de la posición de la curva de demanda y excede al costo marginal $c$.
    \item El precio es $P = c$ y $Q^*$ está determinada, pero el tamaño de cada firma es indeterminado.
    \item Solo puede sobrevivir una única empresa en la industria, configurando automáticamente un monopolio natural.
\end{opts}

\q{Tras una caída permanente de la demanda en una industria competitiva de costos crecientes inicialmente en equilibrio de largo plazo, la secuencia correcta de ajuste es:}
\begin{opts}
    \item En CP cae $P$ con pérdidas; luego salen firmas, sube $P$ y bajan insumos hasta un $P_{LP}$ menor.
    \item En CP sube $P$ por la menor escala; luego entran firmas y el precio final queda por encima del nivel inicial.
    \item En CP cae $P$ y las firmas obtienen beneficios nulos; en LP el precio retorna exactamente al nivel original.
    \item En CP la oferta se contrae instantáneamente y el precio no sufre ninguna variación durante el ajuste.
\end{opts}

\q{¿Cuál es la relación exacta entre el excedente del productor ($EP$) y el beneficio económico ($\pi$) de una empresa competitiva en el corto y en el largo plazo?}
\begin{opts}
    \item En ambos plazos $EP$ es estrictamente menor que $\pi$ porque no descuenta los costos variables totales.
    \item En ambos plazos $EP = \pi$, dado que los costos fijos se recuperan íntegramente mediante el precio.
    \item En corto plazo $EP = \pi - CF$, mientras que en largo plazo $EP$ es igual a los costos hundidos.
    \item En corto plazo $EP = \pi + CF$, mientras que en largo plazo sin costos fijos coincide con $\pi$.
\end{opts}

\q{Un equilibrio competitivo de largo plazo cumple simultáneamente con la eficiencia asignativa y la eficiencia productiva porque:}
\begin{opts}
    \item El precio iguala al costo fijo medio y todas las firmas operan en el tramo creciente del $CMe_{LP}$.
    \item El precio iguala al costo marginal ($P = CMg$) y cada firma produce en el mínimo del $CMe_{LP}$.
    \item La tasa marginal de sustitución supera al precio relativo y los beneficios contables son iguales a cero.
    \item El excedente del consumidor iguala exactamente al excedente del productor en el punto de equilibrio.
\end{opts}

\q{De acuerdo con el Lema de Hotelling aplicado a la función de beneficios $\pi(P, w, r)$ de una empresa perfectamente competitiva, la derivada parcial $\frac{\partial \pi}{\partial P}$ es igual a:}
\begin{opts}
    \item El costo marginal de la empresa evaluado en la escala mínima eficiente de largo plazo.
    \item El negativo de la función de demanda condicionada del factor trabajo en el corto plazo.
    \item La función de oferta óptima $q^*(P, w, r)$ que maximiza el beneficio de la empresa.
    \item La elasticidad precio de la oferta multiplicada por el costo variable medio de la firma.
\end{opts}

\q{La función de beneficios $\pi(P, w, r)$ de una empresa competitiva es convexa en el vector de precios $(P, w, r)$, lo cual implica que:}
\begin{opts}
    \item La oferta óptima $q^*$ nunca disminuye cuando sube el precio del bien ($\partial q^*/\partial P \ge 0$).
    \item Los beneficios siempre caen cuando todos los precios de productos e insumos se duplican simultáneamente.
    \item La empresa prefiere un precio estable al promedio de dos precios fluctuantes con igual media.
    \item La función de costos totales es necesariamente cóncava respecto de la cantidad producida $q$.
\end{opts}

\q{Si tanto el precio del bien final $P$ como todos los precios de los insumos $(w, r)$ se multiplican por un escalar $\lambda > 0$, la función de beneficios $\pi(P, w, r)$ y la oferta $q^*(P, w, r)$ son, respectivamente:}
\begin{opts}
    \item Homogénea de grado cero en precios y homogénea de grado uno en el nivel de producción.
    \item Homogénea de grado uno en precios y homogénea de grado uno en precios de los factores.
    \item Homogénea de grado dos en precios y homogénea de grado menos uno en precios relativos.
    \item Homogénea de grado uno en precios y homogénea de grado cero en precios.
\end{opts}

\q{El área bajo la curva de demanda marshalliana y sobre el precio de mercado ($EC$) coincide exactamente con la variación compensatoria y la variación equivalente únicamente cuando:}
\begin{opts}
    \item La oferta del mercado es perfectamente elástica en el largo plazo y no existen costos hundidos.
    \item El efecto ingreso sobre la demanda del bien es nulo (preferencias cuasilineales).
    \item El bien analizado es un bien Giffen cuya curva de demanda tiene pendiente positiva en todo tramo.
    \item La elasticidad precio de la demanda es constante e igual a menos uno en todos los precios.
\end{opts}

\q{Afirmar que el equilibrio perfectamente competitivo maximiza el bienestar social $W = EC + EP$ significa que:}
\begin{opts}
    \item La distribución resultante del ingreso es justa según cualquier función de bienestar social rawlsiana.
    \item Tanto consumidores como productores reciben exactamente la mitad del excedente total generado.
    \item No existe otro volumen de producción que eleve la suma de excedentes en el mercado.
    \item Todos los consumidores que desean el bien logran adquirirlo independientemente de su disposición a pagar.
\end{opts}

\q{En un mercado competitivo con dos firmas que tienen costos marginales crecientes distintos, la minimización del costo total de producir una cantidad agregada $Q$ se cumple en equilibrio porque:}
\begin{opts}
    \item Ambas enfrentan el mismo precio $P$, igualando sus costos marginales ($CMg_1 = CMg_2 = P$).
    \item La firma con menor costo fijo absorbe toda la producción del mercado y la otra cierra automáticamente.
    \item Ambas firmas producen exactamente la misma cantidad física independientemente de sus costos marginales.
    \item El costo medio de la firma más pequeña iguala al costo variable total de la firma más grande.
\end{opts}

\q{En un mercado perfectamente competitivo, trasladar la obligación legal de pagar un impuesto específico $t$ de los productores a los consumidores:}
\begin{opts}
    \item Aumenta la pérdida irrecuperable de eficiencia porque la demanda se desplaza hacia abajo.
    \item Transfiere íntegramente la carga económica hacia los consumidores sin importar las elasticidades.
    \item Reduce la recaudación fiscal total debido a que los consumidores evaden más fácilmente el tributo.
    \item Deja inalterados los precios netos $P_c$ y $P_p$, la cantidad transada, la recaudación y la $PIE$.
\end{opts}

\q{Para un impuesto específico pequeño $t$ en un mercado competitivo con elasticidad demanda $\varepsilon_d < 0$ y oferta $\varepsilon_s > 0$, la fracción de la carga tributaria soportada por los consumidores es:}
\begin{opts}
    \item $\frac{|\varepsilon_d|}{|\varepsilon_d| + \varepsilon_s}$, por lo que los consumidores pagan más cuando su demanda es más elástica.
    \item $\frac{\varepsilon_s}{\varepsilon_s + |\varepsilon_d|}$, por lo que el lado más inelástico absorbe la mayor parte del impuesto.
    \item $\frac{\varepsilon_s}{|\varepsilon_d|}$, de modo que la carga del consumidor puede superar el 100\% del impuesto en competencia.
    \item $\frac{1}{\varepsilon_s - |\varepsilon_d|}$, siempre que la oferta sea más inelástica que la demanda de mercado.
\end{opts}

\q{En una industria competitiva de costos constantes en el largo plazo (oferta perfectamente elástica), la imposición de un impuesto específico $t$ por unidad implica que:}
\begin{opts}
    \item Los productores absorben el 100\% del impuesto porque no pueden trasladarlo al precio de equilibrio.
    \item La carga del impuesto se reparte en partes exactamente iguales entre consumidores y productores.
    \item Los consumidores absorben el 100\% ($\Delta P_c = t$) y existe $PIE > 0$ si $|\varepsilon_d| > 0$.
    \item La pérdida irrecuperable de eficiencia es nula porque el precio recibido por los productores no cambia.
\end{opts}

\q{Si un bien competitivo tiene una curva de demanda perfectamente inelástica ($\varepsilon_d = 0$) y una oferta con pendiente positiva, un impuesto específico $t$ genera:}
\begin{opts}
    \item Traslación total del impuesto al consumidor ($\Delta P_c = t$) y pérdida irrecuperable de eficiencia nula.
    \item Traslación total del impuesto al productor ($\Delta P_p = -t$) y pérdida irrecuperable de eficiencia máxima.
    \item Una reducción proporcional de la cantidad transada acompañada de una caída en la recaudación fiscal.
    \item Un reparto equitativo del impuesto entre compradores y vendedores con una $PIE$ positiva.
\end{opts}

\q{De acuerdo con la aproximación del Triángulo de Harberger, si se triplica la tasa de un impuesto específico $t$ en un mercado competitivo lineal, la pérdida irrecuperable de eficiencia ($PIE$):}
\begin{opts}
    \item Se triplica exactamente, manteniendo una relación estrictamente lineal con la tasa impositiva.
    \item Se multiplica por seis, porque la recaudación crece en forma cuadrática con la cantidad.
    \item Se reduce a un tercio, debido a que la base imponible disminuye al subir el impuesto.
    \item Se multiplica por nueve, pues la $PIE$ crece proporcionalmente al cuadrado de la tasa ($t^2$).
\end{opts}

\q{Para minimizar el exceso de gravamen ($PIE$) al recaudar un monto fiscal dado mediante impuestos sobre bienes competitivos independientes, la teoría de tributación óptima sugiere:}
\begin{opts}
    \item Gravar con tasas más altas a los bienes con demanda más elástica para reducir el consumo superfluo.
    \item Gravar con tasas inversamente proporcionales a la elasticidad precio de la demanda (Regla de Ramsey).
    \item Aplicar siempre una tasa específica idéntica en unidades monetarias a todos los bienes de la economía.
    \item Gravar únicamente los bienes cuya curva de oferta sea perfectamente elástica en el corto plazo.
\end{opts}

\q{Al comparar un impuesto específico ($t$) y un impuesto ad-valorem ($\tau$) que generan exactamente la misma recaudación en un mercado perfectamente competitivo:}
\begin{opts}
    \item El impuesto ad-valorem genera siempre una pérdida irrecuperable de eficiencia menor que el específico.
    \item El impuesto específico genera siempre un precio al consumidor más bajo y una cantidad mayor.
    \item Ambos impuestos producen exactamente el mismo equilibrio ($Q$, $P_c$, $P_p$) y la misma $PIE$.
    \item El impuesto ad-valorem traslada toda la carga al productor mientras el específico la traslada al consumidor.
\end{opts}

\q{¿Por qué un subsidio específico $s$ por unidad en un mercado competitivo sin externalidades genera una pérdida irrecuperable de eficiencia ($PIE > 0$)?}
\begin{opts}
    \item Porque induce sobreproducción en unidades cuyo costo marginal supera la disposición marginal a pagar.
    \item Porque reduce simultáneamente tanto el excedente del consumidor como el excedente del productor.
    \item Porque desplaza el precio pagado por los consumidores por encima del precio recibido por los productores.
    \item Porque obliga a las empresas competitivas a operar en el tramo decreciente de su costo marginal.
\end{opts}

\q{En un mercado competitivo al que se otorga un subsidio específico $s$ por unidad, ¿qué lado del mercado captura la mayor fracción del beneficio económico del subsidio?}
\begin{opts}
    \item El lado más elástico, porque puede ajustar más rápidamente su cantidad ante el cambio en el precio.
    \item Necesariamente el lado que recibe legalmente el desembolso del Estado en su cuenta bancaria.
    \item Ambos lados capturan exactamente el 50\% del subsidio cuando las curvas son lineales.
    \item El lado más inelástico, exactamente por la misma simetría de elasticidades que en un impuesto.
\end{opts}

\q{En un mercado competitivo con curvas lineales de demanda y oferta, al incrementar progresivamente un impuesto específico $t$ desde cero hasta el nivel prohibitivo, la recaudación fiscal $R(t) = t \cdot Q(t)$:}
\begin{opts}
    \item Crece indefinidamente de manera proporcional al cuadrado de la tasa impositiva $t$.
    \item Aumenta hasta un máximo y luego cae hasta anularse cuando el tributo elimina el mercado.
    \item Permanece constante una vez que la pérdida irrecuperable de eficiencia supera al excedente del productor.
    \item Disminuye desde el primer sol de impuesto siempre que la oferta sea más inelástica que la demanda.
\end{opts}

\q{En una industria competitiva de costos decrecientes ($S_{LP}$ con pendiente negativa), la aplicación de un impuesto específico $t$ puede provocar en el largo plazo que:}
\begin{opts}
    \item El precio al consumidor suba en una cuantía mayor que el impuesto ($\Delta P_c > t$).
    \item El precio pagado por los consumidores disminuya por debajo del precio previo al impuesto.
    \item La cantidad de equilibrio de largo plazo aumente debido a la entrada de nuevas empresas.
    \item La pérdida irrecuperable de eficiencia se vuelva negativa y aumente el bienestar total.
\end{opts}

\q{Cuando el gobierno fija un precio máximo efectivo ($P_{\max} < P^*$) en un mercado perfectamente competitivo, el efecto sobre el excedente del consumidor ($EC$) es:}
\begin{opts}
    \item Siempre positivo, porque todos los consumidores logran comprar a un precio más bajo que $P^*$.
    \item Siempre negativo, porque la escasez elimina por completo las ganancias del intercambio.
    \item Ambiguo: depende de si la ganancia por menor precio supera la pérdida por menor cantidad.
    \item Exactamente nulo, ya que la transferencia desde los productores compensa exactamente la $PIE$.
\end{opts}

\q{Si un precio máximo efectivo genera escasez y el bien se asigna mediante colas (tiempo de espera) en lugar de ir a los consumidores con mayor disposición a pagar, entonces:}
\begin{opts}
    \item La pérdida irrecuperable de eficiencia es menor porque el tiempo de espera sustituye al precio monetario.
    \item La pérdida irrecuperable de eficiencia es estrictamente mayor que el triángulo clásico de Harberger.
    \item El excedente del productor aumenta debido a que las colas incrementan la demanda efectiva del mercado.
    \item El mercado alcanza un nuevo equilibrio eficiente sin necesidad de levantar el control de precios.
\end{opts}

\q{En un mercado competitivo con un precio mínimo efectivo ($P_{\min} > P^*$) sin compras estatales del excedente, la cantidad efectivamente transada en el mercado queda determinada por:}
\begin{opts}
    \item La curva de oferta evaluada en $P_{\min}$, ya que los productores siempre venden todo lo que producen.
    \item El promedio aritmético entre la cantidad demandada y la cantidad ofrecida al precio mínimo.
    \item La intersección entre el precio mínimo y la curva de costo medio mínimo de largo plazo.
    \item La curva de demanda evaluada en $P_{\min}$ (el lado corto del mercado), generando exceso de oferta.
\end{opts}

\q{Si el regulador impone una cuota máxima de producción $Q_c < Q^*$ en un mercado competitivo y subasta los permisos de producción competitivamente, entonces:}
\begin{opts}
    \item El precio sube a $P^d(Q_c)$, el permiso vale $P^d(Q_c) - P^s(Q_c)$ y se genera $PIE$.
    \item El precio de mercado cae hasta $P^s(Q_c)$ y los consumidores capturan íntegramente la renta de la cuota.
    \item La pérdida irrecuperable de eficiencia se elimina porque la subasta transfiere la renta al Estado.
    \item El excedente del productor se reduce a cero debido a que el precio del permiso iguala al costo medio.
\end{opts}

\q{Cuando un país pequeño abre al comercio exterior un mercado competitivo cuyo precio en autarquía supera al precio mundial ($P_{\text{aut}} > P_w$), entonces:}
\begin{opts}
    \item El país exporta el bien, aumenta el excedente del productor y cae el excedente del consumidor.
    \item El país importa el bien; sube $EC$, cae $EP$ y el bienestar total neto aumenta estrictamente.
    \item Tanto el excedente del consumidor como el del productor aumentan en la misma proporción.
    \item El bienestar total disminuye porque la producción nacional se contrae frente a las importaciones.
\end{opts}

\q{En un país pequeño importador, la imposición de un arancel específico $t$ genera dos triángulos de pérdida irrecuperable de eficiencia que representan, respectivamente:}
\begin{opts}
    \item La caída en la recaudación tributaria y el incremento en los costos de transporte internacional.
    \item La pérdida de beneficios de los exportadores extranjeros y la caída del tipo de cambio real.
    \item La distorsión productiva (sobreproducción local) y la distorsión del consumo reducido.
    \item El costo administrativo de aduanas y la pérdida de excedente del productor nacional.
\end{opts}

\q{En competencia perfecta, un arancel específico y una cuota de importación que reducen las importaciones al mismo volumen producen igual precio, consumo y producción local. Su principal diferencia radica en que:}
\begin{opts}
    \item La cuota genera una pérdida irrecuperable de eficiencia dos veces mayor que la del arancel.
    \item El arancel siempre reduce el excedente del productor nacional mientras la cuota lo incrementa.
    \item La cuota impide que el precio interno supere al precio mundial cuando aumenta la demanda local.
    \item Con arancel el rectángulo $(P_d - P_w)M$ es recaudación; con cuota gratuita es renta privada.
\end{opts}

\q{A diferencia de un país pequeño, un país grande con poder de mercado en el comercio internacional puede elevar su bienestar nacional mediante un arancel positivo pequeño si:}
\begin{opts}
    \item La mejora en los términos de intercambio (caída de $P_w$) supera las dos pérdidas triangulares.
    \item El arancel elimina completamente las importaciones y retorna la economía a la situación de autarquía.
    \item La oferta de exportaciones del resto del mundo es perfectamente elástica al precio mundial vigente.
    \item La demanda interna del bien importado es perfectamente elástica respecto del precio internacional.
\end{opts}

\q{En el modelo dinámico de la telaraña (cobweb) con oferta rezagada $Q_t^s = S(P_{t-1})$ y demanda instantánea $Q_t^d = D(P_t)$ lineales, el precio oscila y converge al equilibrio de largo plazo si y solo si:}
\begin{opts}
    \item La curva de demanda es más empinada (más inelástica) en valor absoluto que la curva de oferta.
    \item La curva de oferta es más empinada en valor absoluto que la demanda ($|dP/dQ_s| > |dP/dQ_d|$).
    \item Las pendientes de la oferta y de la demanda son exactamente iguales en valor absoluto.
    \item Los productores forman expectativas racionales perfectas sobre el precio de equilibrio futuro.
\end{opts}

\q{En el modelo de la telaraña con funciones lineales $Q_t^d = a - \beta P_t$ y $Q_t^s = c + \delta P_{t-1}$, si se cumple exactamente $\delta = \beta$, la trayectoria temporal de precios exhibe:}
\begin{opts}
    \item Convergencia monótona sin oscilaciones hacia el precio de equilibrio de largo plazo.
    \item Oscilaciones explosivas de amplitud creciente que alejan indefinidamente al mercado del equilibrio.
    \item Un ciclo perpetuo de amplitud constante alrededor del equilibrio, sin converger ni divergir.
    \item Estabilidad instantánea en un solo período cualquiera sea el precio inicial $P_0$.
\end{opts}

\q{Aunque la distribución de agua potable por redes constituye un monopolio natural, ¿por qué el modelo de competencia perfecta es un referente central para el regulador (SUNASS)?}
\begin{opts}
    \item Porque permite dividir físicamente las redes de distribución en decenas de empresas rivales por manzana.
    \item Porque fija el estándar de eficiencia ($P = CMg$ y $\min CMe$) para evaluar pérdidas de bienestar.
    \item Porque demuestra que las empresas prestadoras de saneamiento nunca requieren inversiones con costos hundidos.
    \item Porque justifica eliminar cualquier forma de regulación tarifaria dejando los precios al libre mercado.
\end{opts}

\q{El mecanismo de competencia por comparación (yardstick competition) utilizado en regulación de servicios de saneamiento consiste en:}
\begin{opts}
    \item Fijar la tarifa de cada EPS según los costos de pares comparables, simulando presión competitiva.
    \item Permitir que los usuarios cambien semanalmente de empresa proveedora de alcantarillado en su vivienda.
    \item Reembolsar automáticamente el 100\% de los costos contables incurridos por cada EPS sin auditoría.
    \item Establecer una tarifa única nacional idéntica sin considerar diferencias geográficas ni topográficas.
\end{opts}

\q{En la regulación tarifaria de saneamiento basada en el enfoque de empresa modelo eficiente, el regulador busca replicar una propiedad clave del equilibrio competitivo de largo plazo:}
\begin{opts}
    \item Que las empresas prestadoras obtengan rentas monopólicas crecientes para financiar gastos corrientes.
    \item Que la demanda de agua potable se vuelva perfectamente elástica en todos los bloques de consumo.
    \item Que las tarifas cubran únicamente los costos operativos variables excluyendo toda inversión en infraestructura.
    \item Que la tarifa remunere solo los costos eficientes y el retorno normal al capital invertido.
\end{opts}

\q{Desde el análisis de excedentes en un mercado competitivo de agua potable sin externalidades, subsidiar directamente el precio por metro cúbico ($P < CMg$) frente a otorgar un bono de suma alzada a hogares vulnerables:}
\begin{opts}
    \item Elimina la pérdida irrecuperable de eficiencia porque induce a los usuarios a ahorrar agua potable.
    \item Resulta preferible en eficiencia asignativa porque iguala la disposición marginal a pagar con el costo marginal.
    \item Genera pérdida de eficiencia por sobreconsumo, mientras el bono preserva la señal $P = CMg$.
    \item Reduce el gasto fiscal del Estado manteniendo exactamente el mismo nivel de bienestar del hogar beneficiario.
\end{opts}

\q{En el sector saneamiento, cuando la competencia en el mercado no es viable por subaditividad de costos y altos costos hundidos, la alternativa de competencia por el mercado (Demsetz) consiste en:}
\begin{opts}
    \item Subastar la concesión al postor que ofrezca la menor tarifa o subsidio bajo metas de calidad.
    \item Duplicar las redes troncales de agua potable para fomentar guerras de precios diarias entre operadores.
    \item Aplicar un precio máximo inferior al costo variable medio para eliminar las rentas extraordinarias.
    \item Prohibir la participación privada y financiar el servicio exclusivamente con impuestos ad-valorem.
\end{opts}

\q{En zonas periurbanas sin red pública abastecidas por numerosos camiones cisterna competitivos, si el precio del combustible sube bruscamente, el modelo de competencia perfecta predice en el corto plazo:}
\begin{opts}
    \item Que el precio del agua en cisterna no cambiará porque el combustible es un costo fijo hundido.
    \item Un alza del $CMg$ y contracción de $S_{CP}$, elevando el precio y reduciendo la cantidad.
    \item Un aumento inmediato en la cantidad ofrecida de agua porque los transportistas necesitan cubrir mayores costos.
    \item Una caída del precio de equilibrio acompañada de un excedente permanente de agua en la zona.
\end{opts}

\q{En una cuenca donde numerosos agricultores competitivos intercambian derechos de agua bien definidos y sin costos de transacción ni externalidades, el equilibrio competitivo del mercado de agua logra:}
\begin{opts}
    \item Que todos los agricultores reciban exactamente el mismo volumen físico de agua por hectárea cultivada.
    \item Que el precio del derecho de agua sea igual a cero incluso en períodos de sequía severa.
    \item Que los agricultores con menor productividad marginal acaparen toda la dotación disponible de la cuenca.
    \item Igualar el valor del producto marginal del agua entre usuarios, maximizando el valor total.
\end{opts}

\subsection*{Parte II. Aplicaciones prácticas y ejercicios cuantitativos (154 a 180)}

\q{Una empresa competitiva tiene la función de costos de corto plazo $C(q) = \frac{1}{3}q^3 - 5q^2 + 24q + 40$. Si el precio de mercado es $P = 15$, ¿cuál es la cantidad óptima que maximiza su beneficio?}
\begin{opts}
    \item $q^* = 1$ unidad
    \item $q^* = 5$ unidades
    \item $q^* = 9$ unidades
    \item $q^* = 10$ unidades
\end{opts}

\q{Dada la función de costo total de corto plazo $C(q) = q^3 - 12q^2 + 50q + 100$ de una empresa perfectamente competitiva, ¿cuál es su precio de cierre en el corto plazo?}
\begin{opts}
    \item $P_{\text{cierre}} = 14$
    \item $P_{\text{cierre}} = 18$
    \item $P_{\text{cierre}} = 26$
    \item $P_{\text{cierre}} = 50$
\end{opts}

\q{Una empresa competitiva enfrenta la función de costos $C(q) = 0.5q^2 + 8q + 72$. ¿A qué precio de mercado la empresa alcanza su punto de nivelación (beneficio económico exactamente nulo)?}
\begin{opts}
    \item $P = 12.00$
    \item $P = 16.00$
    \item $P = 24.00$
    \item $P = 20.00$
\end{opts}

\q{Una industria competitiva tiene 20 empresas tipo A con $CMg_A(q) = 2q + 10$ (cierre en $P=10$) y 30 empresas tipo B con $CMg_B(q) = 3q + 15$ (cierre en $P=15$). Si $P = 25$, la cantidad total ofrecida por la industria es:}
\begin{opts}
    \item $Q^s = 150$ unidades
    \item $Q^s = 250$ unidades
    \item $Q^s = 225$ unidades
    \item $Q^s = 300$ unidades
\end{opts}

\q{En una industria competitiva operan $N = 40$ empresas idénticas con $C(q) = 2q^2 + 4q + 18$, y la demanda de mercado es $Q^d = 440 - 10P$. En el equilibrio de corto plazo, el precio y el beneficio económico de cada empresa son:}
\begin{opts}
    \item $P^* = 20$ y $\pi^* = 14$
    \item $P^* = 24$ y $\pi^* = 50$
    \item $P^* = 24$ y $\pi^* = 32$
    \item $P^* = 28$ y $\pi^* = 32$
\end{opts}

\q{En una industria competitiva con libre entrada y costos constantes, cada firma tiene $C_{LP}(q) = q^3 - 8q^2 + 26q$ y la demanda de mercado es $Q^d = 500 - 20P$. En el equilibrio de largo plazo, el precio y el número de empresas son:}
\begin{opts}
    \item $P_{LP} = 10$ y $N^* = 75$
    \item $P_{LP} = 10$ y $N^* = 60$
    \item $P_{LP} = 12$ y $N^* = 65$
    \item $P_{LP} = 14$ y $N^* = 55$
\end{opts}

\q{Una industria competitiva de costos constantes tiene empresas con $C_{LP}(q) = q^3 - 8q^2 + 26q$. Si la demanda de mercado se expande de $Q_1^d = 500 - 20P$ a $Q_2^d = 660 - 20P$, ¿cuántas empresas adicionales ingresan en el largo plazo?}
\begin{opts}
    \item $20$ empresas adicionales
    \item $30$ empresas adicionales
    \item $50$ empresas adicionales
    \item $40$ empresas adicionales
\end{opts}

\q{Cada firma de una industria competitiva de costos constantes tiene $C_{LP}(q) = q^2 + 4q + 36$. Si el Estado establece una licencia fija anual de $T = 28$ por firma, en el nuevo equilibrio de largo plazo el precio y la producción por empresa son:}
\begin{opts}
    \item $P_{LP} = 16$ y $q_{EME} = 6$
    \item $P_{LP} = 20$ y $q_{EME} = 8$
    \item $P_{LP} = 20$ y $q_{EME} = 6$
    \item $P_{LP} = 24$ y $q_{EME} = 8$
\end{opts}

\q{En un equilibrio competitivo de largo plazo donde el precio determinado por las empresas marginales es $P_{LP} = 14$, una empresa inframarginal más eficiente tiene costos $C(q) = q^2 + 2q + 16$. ¿Cuál es su renta económica?}
\begin{opts}
    \item $12$ unidades monetarias
    \item $16$ unidades monetarias
    \item $20$ unidades monetarias
    \item $28$ unidades monetarias
\end{opts}

\q{Una empresa competitiva tiene $C(q) = 2q^2 + 6q + 90$ en el corto plazo y enfrenta un precio de mercado $P = 26$. ¿Cuál es el valor de su cuasi-renta (excedente del productor)?}
\begin{opts}
    \item $EP = 50$
    \item $EP = 40$
    \item $EP = 90$
    \item $EP = 130$
\end{opts}

\q{Un mercado competitivo presenta demanda inversa $P^d = 100 - 2Q$ y oferta inversa $P^s = 10 + 3Q$. En equilibrio, el excedente del consumidor ($EC$) y el bienestar social total ($W = EC + EP$) son:}
\begin{opts}
    \item $EC = 324$ y $W = 648$
    \item $EC = 486$ y $W = 810$
    \item $EC = 360$ y $W = 900$
    \item $EC = 324$ y $W = 810$
\end{opts}

\q{En un mercado competitivo con $Q^d = 200 - 4P$ y $Q^s = 6P - 50$, se establece un impuesto específico de $t = 15$ por unidad. En el nuevo equilibrio, el precio al consumidor ($P_c$), al productor ($P_p$) y la cantidad ($Q_t$) son:}
\begin{opts}
    \item $P_c = 31$, $P_p = 16$ y $Q_t = 76$
    \item $P_c = 34$, $P_p = 19$ y $Q_t = 64$
    \item $P_c = 35$, $P_p = 20$ y $Q_t = 60$
    \item $P_c = 40$, $P_p = 25$ y $Q_t = 40$
\end{opts}

\q{Considere un mercado competitivo con demanda $P^d = 80 - Q$ y oferta $P^s = 20 + 2Q$. Si se aplica un impuesto específico de $t = 12$ por unidad, la recaudación fiscal ($R$) y la pérdida irrecuperable de eficiencia ($PIE$) son:}
\begin{opts}
    \item $R = 160$ y $PIE = 36$
    \item $R = 240$ y $PIE = 24$
    \item $R = 192$ y $PIE = 24$
    \item $R = 192$ y $PIE = 48$
\end{opts}

\q{En un mercado competitivo donde la elasticidad precio de la demanda es $\varepsilon_d = -0.6$ y la elasticidad de la oferta es $\varepsilon_s = 1.8$, se aplica un impuesto específico de $t = S/\ 8.00$ por unidad. ¿Cuánto aumenta aproximadamente el precio al consumidor?}
\begin{opts}
    \item $S/\ 6.00$ por unidad
    \item $S/\ 2.00$ por unidad
    \item $S/\ 4.00$ por unidad
    \item $S/\ 5.50$ por unidad
\end{opts}

\q{En un mercado competitivo con $Q^d = 168 - 3P_c$ y $Q^s = 2P_p$, se introduce un impuesto ad-valorem de $\tau = 20\%$ sobre el precio recibido por el productor ($P_c = 1.2 P_p$). En equilibrio, los precios $P_p$ y $P_c$ son:}
\begin{opts}
    \item $P_p = 28$ y $P_c = 33.6$
    \item $P_p = 32$ y $P_c = 38.4$
    \item $P_p = 25$ y $P_c = 30.0$
    \item $P_p = 30$ y $P_c = 36.0$
\end{opts}

\q{En un mercado competitivo con curvas lineales de oferta y demanda, un impuesto específico de $t_1 = 4$ genera una pérdida irrecuperable de eficiencia $PIE_1 = 16$. Si el impuesto se eleva a $t_2 = 10$, la nueva $PIE_2$ será:}
\begin{opts}
    \item $PIE_2 = 40$
    \item $PIE_2 = 100$
    \item $PIE_2 = 64$
    \item $PIE_2 = 160$
\end{opts}

\q{Un mercado competitivo tiene $P^d = 60 - 0.5Q$ y $P^s = 10 + 0.5Q$. Si el Estado otorga un subsidio específico de $s = 10$ por unidad, el gasto fiscal total ($G$) y la pérdida irrecuperable de eficiencia ($PIE$) son:}
\begin{opts}
    \item $G = 500$ y $PIE = 25$
    \item $G = 600$ y $PIE = 100$
    \item $G = 600$ y $PIE = 50$
    \item $G = 700$ y $PIE = 50$
\end{opts}

\q{En un mercado competitivo con $Q^d = 150 - 3P$ y $Q^s = 2P$, se fija un precio máximo de $P_{\max} = 24$. Bajo racionamiento eficiente, la escasez resultante y la pérdida irrecuperable de eficiencia ($PIE$) son:}
\begin{opts}
    \item Escasez $= 30$ y $PIE = 60$
    \item Escasez $= 30$ y $PIE = 90$
    \item Escasez $= 18$ y $PIE = 60$
    \item Escasez $= 24$ y $PIE = 48$
\end{opts}

\q{Considere $Q^d = 140 - 2P$ y $Q^s = 3P - 60$. Si el Estado fija un precio mínimo $P_{\min} = 48$ y se compromete a comprar todo el excedente de oferta a ese precio, el exceso de oferta y el costo fiscal son:}
\begin{opts}
    \item Exceso $= 32$ y costo fiscal $= 1536$
    \item Exceso $= 40$ y costo fiscal $= 1600$
    \item Exceso $= 48$ y costo fiscal $= 2304$
    \item Exceso $= 40$ y costo fiscal $= 1920$
\end{opts}

\q{En un país pequeño, el mercado interno tiene $Q^d = 120 - 2P$ y $Q^s = 2P - 20$. Si el precio mundial es $P_w = 20$, al abrirse al libre comercio el volumen de importaciones ($M$) y la ganancia neta de bienestar frente a la autarquía son:}
\begin{opts}
    \item $M = 50$ y $\Delta W = 375$
    \item $M = 60$ y $\Delta W = 450$
    \item $M = 60$ y $\Delta W = 900$
    \item $M = 80$ y $\Delta W = 600$
\end{opts}

\q{Un país pequeño enfrenta $Q^d = 120 - 2P$, $Q^s = 2P - 20$ y precio mundial $P_w = 20$. Si establece un arancel específico de $t = 5$ por unidad importada, la recaudación arancelaria ($R$) y la pérdida irrecuperable de eficiencia ($PIE$) son:}
\begin{opts}
    \item $R = 200$ y $PIE = 25$
    \item $R = 300$ y $PIE = 50$
    \item $R = 200$ y $PIE = 50$
    \item $R = 250$ y $PIE = 75$
\end{opts}

\q{En una economía pequeña con $Q^d = 120 - 2P$, $Q^s = 2P - 20$ y $P_w = 20$, el gobierno fija una cuota máxima de importación de $M = 20$ unidades. El precio interno resultante y la renta de cuota son:}
\begin{opts}
    \item $P = 30$ y renta de cuota $= 200$
    \item $P = 30$ y renta de cuota $= 300$
    \item $P = 25$ y renta de cuota $= 100$
    \item $P = 35$ y renta de cuota $= 300$
\end{opts}

\q{En un mercado competitivo tipo telaraña con $Q_t^d = 100 - 2P_t$ y $Q_t^s = 10 + P_{t-1}$, si el precio inicial es $P_0 = 20$, los precios en los períodos $t=1$ y $t=2$ son:}
\begin{opts}
    \item $P_1 = 30.0$ y $P_2 = 30.0$
    \item $P_1 = 35.0$ y $P_2 = 25.0$
    \item $P_1 = 40.0$ y $P_2 = 20.0$
    \item $P_1 = 35.0$ y $P_2 = 27.5$
\end{opts}

\q{La función de beneficios de corto plazo de una firma competitiva es $\pi(P, w) = \frac{P^2}{4w} - 10$. Si $P = 20$ y $w = 5$, utilizando el Lema de Hotelling, la oferta óptima $q^*$ y la demanda de trabajo $L^*$ son:}
\begin{opts}
    \item $q^* = 4$ y $L^* = 2$
    \item $q^* = 2$ y $L^* = 4$
    \item $q^* = 2$ y $L^* = 2$
    \item $q^* = 4$ y $L^* = 8$
\end{opts}

\q{En un mercado competitivo con oferta perfectamente inelástica $Q^s = 80$ y demanda $Q^d = 200 - 3P$, se impone un impuesto específico de $t = 10$ a los consumidores. En el nuevo equilibrio, $P_c$, $P_p$ y la $PIE$ son:}
\begin{opts}
    \item $P_c = 50$, $P_p = 40$ y $PIE = 0$
    \item $P_c = 45$, $P_p = 35$ y $PIE = 75$
    \item $P_c = 40$, $P_p = 30$ y $PIE = 0$
    \item $P_c = 40$, $P_p = 30$ y $PIE = 400$
\end{opts}

\q{Para estimar la tarifa de referencia competitiva de largo plazo de un servicio local, el regulador modela una empresa eficiente con $C(q) = 0.02q^2 + 4q + 200$. La escala mínima eficiente ($q_{EME}$) y la tarifa de largo plazo ($P_{LP}$) son:}
\begin{opts}
    \item $q_{EME} = 100$ y $P_{LP} = 8$
    \item $q_{EME} = 100$ y $P_{LP} = 6$
    \item $q_{EME} = 50$ y $P_{LP} = 9$
    \item $q_{EME} = 200$ y $P_{LP} = 9$
\end{opts}

\q{En un mercado competitivo de agua con costo marginal constante $CMg = 6$ y demanda inversa $P^d = 20 - 0.2Q$, se aplica un subsidio tarifario de $s = 2$ por metro cúbico ($P_c = 4$). ¿Cuál es la pérdida irrecuperable de eficiencia ($PIE$) frente a la asignación eficiente?}
\begin{opts}
    \item $PIE = 20$ soles
    \item $PIE = 15$ soles
    \item $PIE = 40$ soles
    \item $PIE = 10$ soles
\end{opts}
